\section{Escritura y orden de trazado}\label{sec:escritura}

\subsection{Forma y tamaño}
Para dibujar cada signo uso una cuadrícula de 64 por 64 unidades, con la parte principal
aproximadamente entre 12 y 52. Mantengo las proporciones, el grosor uniforme y las
uniones redondeadas. No se requieren presión variable, colores ni un instrumento
especial. Un bolígrafo normal debe permitir reproducir los rasgos distintivos.

Como punto de partida para practicar, tomo una caja de 9 a 12 mm por signo y
ajusto el tamaño a mi comodidad; no es una medida ergonómica validada. El tamaño
no cambia el significado. Conservo abiertos los huecos y visibles las barras
cortas, especialmente al reducir la escritura.

\subsection{Espaciado y organización}
Dejo aproximadamente un cuarto del ancho de un signo entre signos independientes.
Los índices van abajo a la derecha, cerca de su signo y separados del siguiente
argumento. Evito integrarlos en el contorno del glifo. Dejo margen entre el borde
de un grupo y su contenido; un borde nunca debe confundirse con un trazo del signo.

Hay tres estructuras externas: grupo de borde continuo, región convencional
de borde discontinuo y barra vertical entre argumentos. A mano, uso corchetes
completos para el grupo y conservo un contorno discontinuo para la región
convencional. Los distingo por el borde, con independencia del color.

Una unidad completa ocupa un bloque separado. Los grupos anidados deben quedar
visibles; evito cortar una expresión entre páginas. Si una composición es
incómodamente larga, defino antes una abreviatura con K09. El contexto del
documento continúa en la página siguiente según la gramática.

\subsection{Cómo seguir las instrucciones}
En la tabla reúno las formas vigentes y propongo un orden de trayectos.
Dentro de cada número mantengo el lápiz sobre el papel; entre números lo
levanto. Izquierda, derecha, arriba y abajo se refieren al signo tal como aparece
impreso. Los recuentos describen el dibujo propuesto, no movimientos medidos de
la mano ni una velocidad comprobada. Puede variarse el orden si se conserva la
identidad del signo; todavía no se ha fijado una variante cursiva enlazada.

% Orden orientativo de ejecución; una fila por signo canónico.
\newcommand{\KalaTrazoFila}[2]{%
\KalaSimbolo[13mm]{#1}\quad\texttt{#1} &
\textbf{\csname KalaNombre#1\endcsname}
(\csname KalaTrayectos#1\endcsname\
\ifnum\csname KalaTrayectos#1\endcsname=1 trayecto\else trayectos\fi). #2\\[7pt]}
\begin{longtable}{@{}p{32mm}>{\raggedright\arraybackslash}p{\dimexpr\linewidth-32mm-2\tabcolsep\relax}@{}}
\toprule Modelo & Orden de trazado \\ \midrule\endhead
\bottomrule\endfoot
\KalaTrazoFila{K01}{1. Bajo por la izquierda y curvo hacia la derecha en la base.
2. Dibujo el ángulo superior desde la izquierda hacia la derecha y bajo.}
\KalaTrazoFila{K02}{1. Desde abajo a la izquierda, subo curvando por arriba hacia
la derecha y vuelvo a la izquierda hasta el centro. 2. Bajo desde ese extremo.}
\KalaTrazoFila{K03}{1. Trazo la curva ascendente de izquierda a derecha y termino
horizontal. 2. Dibujo la punta desde arriba a la izquierda, pasando por el
extremo derecho y bajando a la izquierda.}
\KalaTrazoFila{K04}{1. Desde abajo a la izquierda, dibujo una onda con arco
superior, valle y subida derecha. 2. Añado la punta superior de izquierda a derecha.}
\KalaTrazoFila{K05}{1. Bajo por la izquierda y giro a la derecha en la base.
2. Añado la barra superior hacia la derecha. 3. Añado la barra intermedia.}
\KalaTrazoFila{K06}{1. Arco izquierdo abierto, de arriba abajo.
2. Zigzag derecho: derecha, diagonal abajo a la izquierda y derecha.}
\KalaTrazoFila{K07}{1. Desde arriba a la izquierda, curvo por la derecha y vuelvo
a la izquierda por abajo. 2. Trazo la L interior bajando y girando a la derecha.}
\KalaTrazoFila{K08}{1. Cierro el rombo: izquierda, arriba, derecha, abajo,
izquierda. 2. Barra interior hacia la derecha hasta tocar el lado oblicuo,
sin sobresalir.}
\KalaTrazoFila{K09}{1. Corchete izquierdo: derecha, abajo, izquierda.
2. Barra derecha superior. 3. Barra derecha inferior.}
\KalaTrazoFila{K10}{1. Desde abajo a la izquierda, subo, redondeo la esquina,
sigo a la derecha, bajo y vuelvo hacia el interior a la izquierda.}
\KalaTrazoFila{K11}{1. Arco desde arriba a la izquierda por la derecha hasta
abajo. 2. Diagonal desde abajo a la izquierda hacia arriba a la derecha.}
\KalaTrazoFila{K12}{1. V desde arriba a la izquierda al centro inferior y arriba
a la derecha. 2. Bajo una cola desde el vértice inferior.}
\KalaTrazoFila{K13}{1. Desde arriba a la izquierda, voy al centro y vuelvo abajo
a la izquierda. 2. Del centro hacia la derecha y un retorno corto arriba a la izquierda.}
\KalaTrazoFila{K14}{1. U desde arriba a la izquierda, base curva y subida derecha.
2. Barra horizontal entre los lados. 3. Vertical central desde arriba, cruzando la barra.}
\KalaTrazoFila{K15}{1. La misma U. 2. Barra corta interior, separada de ambos
lados. Conservo esa separación y no añado la vertical de K14.}
\KalaTrazoFila{K16}{1. Desde abajo a la izquierda: subo, derecha, esquina
inclinada, bajo y cierro hacia la izquierda. 2. Barra interior hacia la derecha.}
\KalaTrazoFila{K17}{1. Ángulo izquierdo hacia el centro y de vuelta abajo a la
izquierda. 2. Ángulo derecho hacia el centro y de vuelta abajo a la derecha.
Dejo separados los vértices; no formo una X.}
\KalaTrazoFila{K18}{1. Desde arriba a la derecha, voy a la izquierda, bajo con
esquinas curvas y vuelvo a la derecha por abajo. 2. Ángulo interior hacia la
derecha y abajo a la izquierda.}
\KalaTrazoFila{K19}{1. Desde abajo a la izquierda, curvo por arriba hacia la
derecha, vuelvo en diagonal abajo a la izquierda y termino abajo a la derecha.}
\KalaTrazoFila{K20}{1. Barra superior hacia la derecha, diagonal abajo a la
izquierda y bajada vertical. 2. Barra intermedia hacia la derecha.}
\KalaTrazoFila{K21}{1. Escalera: subo, derecha, subo, derecha.
2. Punta de flecha desde arriba a la izquierda, extremo y abajo a la izquierda.}
\KalaTrazoFila{K22}{1. Desde arriba a la derecha, rodeo por arriba, izquierda y
abajo, subo por la derecha y termino con la barra hacia dentro. Todo es un
trayecto continuo; no añado una bajada ni repaso el arco.}
\end{longtable}


\subsection{Rasgos que deben conservarse}
Colección mantiene separados el arco y el zigzag. Fuente conserva el contorno
con esquina inclinada y la barra interior. En Espacio la barra llega al lado
oblicuo sin sobresalir. Mientras tiene un único recorrido continuo de arco a
barra: no lleva un segmento adicional superpuesto.

Universal y Existencia se distinguen por la barra larga y vertical de K14 frente
a la barra corta interior de K15. Esa semejanza visual no convierte un operador
en una intensificación del otro. En Alternativa se mantienen separados los dos
vértices centrales.

No hay una regla global de tres alturas, rotaciones con significado ni un
modificador universal de opuestos. El orden sugerido facilita la copia; su
comodidad y recuerdo deben comprobarse con escritura real. Conservo las formas anteriores
y mis observaciones en el anexo gráfico del manual.
